\section{Self-consistent estimate against grain size}\label{sup:selfconsistent}

Figure~\ref{fig:s-selfconsistent} shows the minimisers of the residual $\scres$ of section~7
of the paper for grain radii $R=2.5$--$500$~\AA, with the exact averaged Eshelby tensor of
appendix~B, with seven and eight terms. The residual was evaluated on a grid in $(\nu,\ell)$
and every local minimum was refined by the Nelder--Mead method: each of the $\NumSCCases$
cases has a single minimum, inside the grid. Both unknowns vary smoothly with $R$
and tend to the closed limit of the paper: $\nu$ to Kr\"oner's value and $\ell$ to $\linf$,
within $\NumSCNearLimitPercent$\,\% from $2R=2$~nm on. The normalised residual
$\scres/|\Phi|$ hardly depends on $R$: the mismatch is not a size effect but the
inability of a one-parameter medium to match the block lengths of the grain. In the
seven-term solutions the omitted condition $B_2'=3$ is missed by up to
$\NumSCSixthSolutionMaxPercent$\,\%, less for larger grains. Grains whose radius is below the
lattice parameter (shaded) are not physical and are shown only to display the trend; there
the minimisers depart most from the limit, which involves only the Voigt estimate of the
grain. The Cu MEAM tensor is nearly singular (section~5 of the paper), but the scheme
does not invert it: moving $D_{16,16}$ and $D_{16,17}$, which fix $m_E$, by the rounding of
their published digits changes $\ell$ by less than $\NumCuMEAMSensitivityEll$~\AA, at $R=5$~\AA\ and in the limit.

\begin{figure}[pos=htbp]
\centering
\includegraphics{selfconsistent_fcc}
\caption{Self-consistent estimate against the grain diameter $2R$: (a) effective length
$\ell$, (b) Poisson's ratio $\nu$ and (c) normalised residual $\scres/|\Phi|$, with seven
terms (solid lines, filled markers) and eight (dashed lines, open markers). Dotted lines:
the closed large-grain limit $\linf$ (seven terms) in (a) and Kr\"oner's $\nu_{\mathrm K}$
in (b). Shaded: radius below the lattice parameter of Al, the largest of the three
($R<a$), not physical.}\label{fig:s-selfconsistent}
\end{figure}
